\section{Related Work}

\paragraph{Execution Traces for Code LLMs.}
TracePile \citep{chen2025tracepile} scaled trace-augmented training to 2.6M samples; \citet{wang2025dcsh} found trace utility task- and format-dependent; relative to DCSH, we contribute (i) a two-tier signal hierarchy quantified via model-free diagnostic token ratio ($r{=}0.949$), (ii)~a phase transition in trace length separating effective from ineffective truncation budgets, (iii)~calibration analysis (ECE from 0.171 to 0.003), and (iv)~cross-dataset and cross-architecture validation across four models.
NExT \citep{ni2024next} and ET-CoT \citep{eguchi2025etcot} used trace chain-of-thought for repair and execution reasoning; \citet{ferrando2025execute} found compact traces comparable to full traces, supporting our density finding.
\citet{liu2023codeexecution} pre-trained on execution traces; CRUXEval \citep{gu2024cruxeval}, extended to 19 languages \citep{xu2025cruxevalx}, established execution reasoning benchmarks.
None decompose which trace components contribute to \emph{classification} accuracy.
\citet{shi2023distracted} showed LLMs are easily distracted by irrelevant context, providing a complementary perspective on why verbose traces---which bury diagnostic tokens in thousands of irrelevant lines---impair SFT learning.

\paragraph{Code Verdict and Patch Classification.}
TRACED \citep{ding2024traced} showed execution-aware pretraining improves vulnerability detection; DeepDFA \citep{steenhoek2024deepdfa} and learned program slicing \citep{yadavally2024slicing} combined classical analysis with deep learning.
In patch correctness, \citet{tian2020patchcorrectness}, CC2Vec \citep{hoang2020cc2vec}, LLM4PatchCorrect \citep{zhou2024llm4patch} (84.4\%), and APPT \citep{appt2024} (79.7\%) developed representation learning of code changes.
Verdict prediction is an instance of the longstanding test oracle problem \citep{barr2015oracle}; FuseFL \citep{widyasari2024fusefl} and \citet{foster2025mutation} demonstrated structured fault signals; the equivalent-mutant problem \citep{papadakis2019mutation} poses an analogous classification challenge.
None isolate individual trace components or compare against rule-based baselines.

\paragraph{Fault Localization.}
Spectrum-based fault localization (SBFL) ranks program elements by suspiciousness from test spectra \citep{abreu2007accuracy,jones2002visualization,wong2014dstar}; recent LLM-based approaches \citep{kang2024autofl,qin2024agentfl} share our insight that anomalous execution events carry the strongest diagnostic signal, with verdict prediction serving as a prerequisite to FL.

\paragraph{Trace Compression and Selection.}
TraceCoder \citep{huang2026tracecoder} and InspectCoder \citep{wang2025inspectcoder} use selective trace extraction; LDB \citep{zhong2024ldb} applies segment-level trace verification for debugging; SelfDebug \citep{chen2024selfdebug} iteratively repairs code with execution feedback; Reflexion \citep{shinn2023reflexion} uses verbal reinforcement from execution outcomes to guide agent learning. RepairAgent \citep{bouzenia2025repairagent} and Agentless \citep{xia2024agentless} incorporate execution feedback in repair workflows, as do SWE-agent \citep{yang2024sweagent} and OpenHands \citep{wang2025openhands}; in the SWE-bench ecosystem \citep{jimenez2024swebench}, these agents rely on test verdict signals to guide patch selection, yet which trace components carry that signal most efficiently remains unstudied.
In the prompt compression literature, LLMLingua \citep{jiang2023llmlingua} and Selective Context \citep{li2023selectivecontext} showed that pruning low-information tokens preserves downstream accuracy, corroborating our finding that compact trace formats outperform verbose ones; \citet{haque2025traces} found raw traces provide limited repair benefit.
Our approximate decomposition (gradient dilution places an upper bound of ${\leq}$4pp on the gap) provides the first quantitative framework for understanding how compression helps, complementing these methods' focus on \emph{how}.

\paragraph{Calibration and Defect Prediction.}
Neural network calibration \citep{guo2017calibration}, extended to LLMs \citep{kadavath2022language,geng2024calibration,spiess2025calibration}, studies confidence-correctness alignment; our ECE ${<}$ 0.007 suggests compact inputs improve calibration. Just-in-time defect prediction \citep{kamei2013large,hoang2019deepjit,pornprasit2021jitline} shares our fault detection motivation but operates on code change features.

\paragraph{Training Efficiency and Data Selection.}
Curriculum learning \citep{bengio2009curriculum}, dataset cartography via training dynamics \citep{swayamdipta2020cartography}, and data selection methods \citep{sachdeva2024dataefficient,xie2023doremi} show that data quality and ordering improve convergence, consistent with our finding that compact exceptions enable rapid convergence while verbose traces trap the model in oscillation.
